%%%%%%%%%%%%%%%%%%%%%%%%%%%%%%%%%%%%%%%%%%%%%%%%%%%%%%%%%%%%%%
%% Rozdział 6: Tłumaczenie na poziomie zdań (video2text)
%%%%%%%%%%%%%%%%%%%%%%%%%%%%%%%%%%%%%%%%%%%%%%%%%%%%%%%%%%%%%%

\section{Rozpoznawanie sekwencji gestów w czasie rzeczywistym}
\label{chap:sentence-level-translation}

Rozdział~\ref{chap:isolated-gloss-recognition} opisuje model, który
przekształca sekwencję punktów charakterystycznych w embedding i wyszukuje
podobne nagrania referencyjne w bazie FAISS. W praktycznej aplikacji dane
pochodzą jednak z nieprzerwanego strumienia kamery. System musi więc
decydować, które fragmenty strumienia przekazywać do modelu i kiedy uznać
predykcję za wystarczająco stabilną. W tym rozdziale opisano dwa
zaimplementowane tryby rozpoznawania w czasie rzeczywistym: segmentację
gestów rozdzielonych pauzą oraz ciągłe rozpoznawanie z wykorzystaniem okna
przesuwnego.

\subsection{Całościowy potok przetwarzania}
\label{sec:sentence-level-pipeline}

Każda klatka strumienia kamery jest analizowana przez MediaPipe Hands i
MediaPipe Pose, a wynik jest zapisywany jako 144-wymiarowy wektor punktów
charakterystycznych. Wybrany fragment sekwencji przechodzi identyczne
przetwarzanie wstępne jak dane treningowe: interpolację braków, usunięcie
klatek statycznych, wygładzanie, normalizację przestrzenną oraz resampling
do 60 klatek. Enkoder neuronowy tworzy następnie 64-wymiarowy embedding,
który jest porównywany z embeddingami referencyjnymi w indeksie FAISS.

W trybie segmentacji rozpoznawanie rozpoczyna się po trzech kolejnych klatkach
z wykrytą dłonią. Punkty są gromadzone w buforze o maksymalnej długości
150 klatek, a segment zostaje zamknięty po dziesięciu kolejnych klatkach bez
wykrytej dłoni albo po zapełnieniu bufora. Sekwencje krótsze niż osiem klatek
są odrzucane. Tryb ten jest przeznaczony dla gestów rozdzielonych wyraźną pauzą.

Tryb ciągły wykorzystuje okno przesuwne. Po wykryciu aktywności dłoni system
rozpoczyna rejestrowanie po jednej aktywnej klatce i utrzymuje bufor pięciu
najnowszych klatek. Po jego zapełnieniu wnioskowanie jest wykonywane co dwie
nowe klatki, dlatego kolejne okna częściowo się pokrywają. Do modelu trafia
zawsze sekwencja po resamplingu do 60 klatek, zgodna z formatem treningowym.
Predykcja jest rozważana tylko wtedy, gdy dłoń wykryto w co najmniej $60\%$
klatek okna.

Wyszukiwanie FAISS zwraca etykietę, podobieństwo \emph{score} oraz margines
\emph{margin} względem najlepszej klasy konkurencyjnej. Dla bazy podstawowej
akceptowane są wyniki o $score \geq 0.90$ i $margin \geq 0.03$; dla bazy
utworzonej z własnych nagrań stosowany jest próg $score \geq 0.75$ przy tym
samym minimalnym marginesie. Aby ograniczyć migotanie wyniku, system analizuje
trzy ostatnie pewne predykcje i akceptuje glosę, gdy co najmniej dwie z nich
wskazują tę samą etykietę, a ostatnia predykcja również ma tę etykietę.
Po zaakceptowaniu glosa jest blokowana przed ponownym dodaniem; zmiana etykiety
wymaga trzech stabilnych predykcji nowej klasy. Dodatkowo stosowany jest
trzyklatkowy okres chłodzenia. Po trzech kolejnych klatkach bez wykrytej dłoni
bufory są czyszczone, a system wraca do stanu oczekiwania.

\subsection{Integracja modułów z rozdziałów~\ref{chap:isolated-gloss-recognition} i~\ref{chap:text-translation}}
\label{sec:sentence-level-integration}

Zaakceptowane glosy są dopisywane do pamięci bieżącej sekwencji i wyświetlane
w interfejsie użytkownika. Użytkownik może wyczyścić całą sekwencję lub usunąć
ostatnią glosę. Pamięć ta stanowi interfejs dla planowanego modułu gloss2text
opisanego w Rozdziale~\ref{chap:text-translation}, lecz aktualna wersja nie
generuje jeszcze zdań języka naturalnego. Zaimplementowano więc ciągłe
rozpoznawanie sekwencji glos, a nie pełny potok video2text.

\subsection{Eksperymenty i wyniki}
\label{sec:sentence-level-experiments}

Implementacja udostępnia tryb diagnostyczny. Dla każdej próby zapisywane są
prawdziwa i przewidziana etykieta, poprawność predykcji, podobieństwo oraz
margines. Pozwala to analizować nie tylko dokładność rozpoznania, lecz również
pewność decyzji podejmowanych przez wyszukiwanie FAISS i stabilizator.

Pełna ocena tłumaczenia video2text wymaga jeszcze modułu gloss2text oraz
nagrań ciągłych z referencyjnymi zdaniami. Po jego dodaniu należy rozdzielić
błędy rozpoznania glos od błędów tłumaczenia sekwencji glos na tekst.